% ─────────────────────────────────────────────────────────────────────────────
% Supplementary: TTFT breakdown (docs/06 §9 #6). Table from results/paper/build_ttft_breakdown.py,
% using the existing efficiency measurements (no extra runs).
% ─────────────────────────────────────────────────────────────────────────────
\section{TTFT breakdown}
\label{sec:supp_ttft}

% 由 results/paper/build_ttft_breakdown.py 產生，不要手改
\begin{table*}[t]
  \centering
  \small
  \setlength{\tabcolsep}{4.5pt}
  \begin{tabular}{lccccccccccc}
    \toprule
    & \multicolumn{5}{c}{(a) Budget $B$ (five benchmarks)} & & \multicolumn{5}{c}{(b) Input tiles (HR-8K)} \\
    \cmidrule(lr){2-6}\cmidrule(lr){8-12}
    Stage (ms) & 1024 & 1536 & 2048 & 2560 & 3072 & & 7 & 25 & 78 & 127 & 188 \\
    \midrule
    Prefill of $\mathbf{x}_b$ & 41 & 41 & 41 & 41 & 41 & & 41 & 41 & 42 & 42 & 43 \\
    Vision encoder, coverage tiles & 20 & 22 & 41 & 43 & 65 & & 22 & 22 & 23 & 24 & 23 \\
    Prior pass & 95 & 101 & 174 & 180 & 246 & & 98 & 99 & 99 & 100 & 99 \\
    Vision encoder, detail tiles & 62 & 70 & 69 & 86 & 88 & & 68 & 70 & 70 & 71 & 72 \\
    Boundary scoring ($q_t$) & 172 & 140 & 149 & 118 & 124 & & 129 & 131 & 137 & 132 & 136 \\
    Prefill of selected tiles & 188 & 290 & 381 & 488 & 580 & & 288 & 291 & 291 & 294 & 292 \\
    Prefill of $\mathbf{x}_a$ + first token & 42 & 42 & 43 & 43 & 44 & & 42 & 42 & 42 & 42 & 42 \\
    \midrule
    \rowcolor{oursbg}
    TTFT, \method & 623 & 709 & 899 & 1002 & 1191 & & 690 & 699 & 706 & 707 & 710 \\
    \quad of which selection (prior + $q_t$) & 43\% & 34\% & 36\% & 30\% & 31\% & & 33\% & 33\% & 33\% & 33\% & 33\% \\
    TTFT, uncompressed & \multicolumn{5}{c}{1740 (independent of $B$)} & & 414 & 1428 & 5258 & 10013 & 17229 \\
    \bottomrule
  \end{tabular}
  \caption{\textbf{TTFT breakdown} of \method on InternVL3.5-8B (mean in ms; the stages add up to the mean TTFT up to about 2\,ms of CPU time). (a) Mean over DocVQA, InfoVQA, MME-RealWorld, HR-Bench 4K, and V$^*$ Bench, 50 images each. (b) HR-Bench 8K with the tile cap raised from 7 to 201, at $B=1536$; columns give the mean number of input tiles including the thumbnail.}
  \label{tab:ttft}
\end{table*}


\cref{tab:ttft} splits the TTFT of \method into the stages of \cref{alg:foveatile}, measured as in
\cref{sec:setup}. The two stages that exist only to select tiles, the prior pass and boundary scoring, take
240 to 270\,ms at $B\le1536$ and 300 to 370\,ms at larger budgets. They replace the encoding and prefill
of the tiles that are never selected.

\paragraph{Budget.} The largest stage is the prefill of the selected tiles, which grows linearly with $B$
(188 to 580\,ms). The prior pass grows with the $m=\lfloor K/2\rfloor$ coverage tiles it reads, and boundary
scoring shrinks as the window $S$ narrows from 5 to 3. The text prefills before and after the image take
about 40\,ms each at every budget.

\paragraph{Resolution.} Raising the tile cap from 7 to 201 changes no stage by more than 10\,ms, and TTFT
stays between 0.69 and 0.71\,s, as \cref{tab:complexity} predicts: every stage depends on $K$ and $S$ and
none on $N$. The detail tiles that fall outside the directly kept set and the window are never encoded,
so the vision encoder runs on the same number of tiles at every resolution. Without compression, TTFT grows
from 1.4\,s at 25 tiles to 17.2\,s at 188 tiles, 24$\times$ the TTFT of \method.
